\section{Performance conditions determine units before settlement}
\label{sec:apple-performance-award}

The performance form turns the distinction between a plan and a grant into a
calculation problem. Its notice leaves the target quantity, grant date,
performance period and vesting schedule blank. Exhibit A likewise leaves the
percentile thresholds and payout percentages blank, although it specifies
interpolation between thresholds, a cap when Apple's absolute return is
negative, and rounding to whole units. The form therefore supplies part of
the method without supplying an individual award's inputs
\citep[Exhibit 10.3, notice and Exhibit A]{apple2022plan}.

The 2025 proxy provides a concrete schedule for the named executives'
1 October 2023 performance awards. It describes a three-year performance
period from the beginning of fiscal 2024 through the end of fiscal 2026, with
vesting on 1 October 2026 subject to the agreements. This is a dated
description of that grant population, not the result of the performance
period. The proxy separately reports results for awards granted on
27 September 2020 and vested on 1 October 2023. Those older results must not
be attached to the newer grants merely because the same proxy reports both
\citep[printed pp.~51--53]{apple2025proxy}.

For the October 2023 schedule, first suppose that the relative-return
percentile has been supplied correctly. Write it as $p$ on a scale from 0 to
100, so that $p=70$ means the seventieth percentile. Let $m(p)$ be the number
of units earned per target unit before applying the negative-return cap and
employment conditions. The proxy gives zero below the twenty-fifth percentile,
25\% of target at that percentile, 100\% at the fifty-fifth and 200\% at the
eighty-fifth or above. Between those points it requires straight-line
interpolation \citep[printed p.~52]{apple2025proxy}. In these units,
\[
 m(p)=\begin{cases}
 0, & p<25,\\[4pt]
 \dfrac14+\dfrac{p-25}{40}, & 25\leq p<55,\\[7pt]
 1+\dfrac{p-55}{30}, & 55\leq p<85,\\[7pt]
 2, & 85\leq p\leq100.
 \end{cases}
\]
The first slope is $(1-1/4)/(55-25)=1/40$; the second is
$(2-1)/(85-55)=1/30$. Thus a supplied percentile of 40 earns 0.625 units per
target unit, while 70 earns 1.5. The threshold at 25 is a jump, not a line
running from zero at some unreported lower percentile. If two unresolved
rank calculations produced only $p=24.9$ or $p=25$, a hypothetical target of
1,000 units would yield either zero or 250 units before other conditions.
Replacing those two possibilities by every quantity from zero to 250 would
create unsupported outcomes.

Relative performance is only one input. Let $R$ be Apple's absolute total
shareholder return over the same performance period, expressed as a fraction.
The source caps vesting at 100\% of target whenever $R<0$, regardless of
relative rank. For this part of the calculation the multiplier is
\[
 g(p,R)=\begin{cases}\min\{m(p),1\},&R<0,\\m(p),&R\geq0.\end{cases}
\]
At the supplied percentile 70, a hypothetical target of 1,000 units therefore
gives 1,500 units when the return is nonnegative but only 1,000 when it is
negative, assuming the other vesting conditions are met. A company can rank
well in a falling market: the rank does not identify the sign of its return.
The example avoids fractional units; Exhibit A's nearest-whole-unit rule still
needs an explicit treatment of exact half-unit ties before implementing all
inputs \citep[Exhibit 10.3, Exhibit A]{apple2022plan}.

Computing $R$ introduces its own source requirements. The form uses
twenty-trading-day beginning and ending average prices, dividend reinvestment,
and adjustments for specified capital changes. Consider a deliberately simple
path with no dividends, distributions or corporate actions. Suppose each of
the twenty beginning prices is 100; nineteen ending-window prices are 80 and
the final price is 120. The beginning average is 100 and the ending average is
$(19\times80+120)/20=82$, giving $R=(82-100)/100=-18\%$. Substituting only
the first and last prices would instead give $+20\%$. With the supplied
percentile held at 70, that substitution would wrongly remove the cap and
increase the hypothetical award from 1,000 to 1,500 units. Even in this simple
case, the observation window changes a contractual consequence
\citep[Exhibit 10.3, Exhibit A, Beginning Price and Ending Price]{apple2022plan}.

The general return and rank calculation requires more than those two averages.
The form defines the comparison group using initial S\&P~500 membership and
continued listing, with specific bankruptcy, successor and merger provisions.
It ranks Apple ahead of a comparison member with equal return. That tie rule
does not by itself specify a unique numerical percentile formula; a library's
default rank conversion cannot silently supply one. Dividends enter the return
through reinvestment, noncash distributions need values, and corporate actions
need the stated adjustments. A calculation must document how those provisions
are implemented, including avoiding a second addition of dividends to a price
series that already includes them. Unresolved ambiguities belong with the
Committee's determination under the form; they are not permission to invent
historical inputs \citep[Exhibit 10.3, Exhibit A]{apple2022plan}.

Employment conditions then determine which otherwise earned units survive.
Under Section 8 of the filed 2022 performance form, ordinary termination
generally cancels unvested rights. Death or disability before vesting instead
allows a fraction of the performance-qualified award to vest on the specified
vesting date. The fraction is the total performance-period days less the days
remaining at severance, divided by total days. For an expressly hypothetical
eligible fraction of $3/5$ and the uncapped 1,500-unit result, 900 units survive
to that date; if the negative-return cap first reduces the result to 1,000,
600 survive. These calculations use the filed form's terms and assume no
other modification, leave adjustment or recoupment event
\citep[Exhibit 10.3, Sections 3--4 and 7--9]{apple2022plan}.

The proxy's termination summary must remain alongside that versioned form,
rather than overwrite it. It describes pro-rata \emph{accelerated} vesting for
named executives other than Cook upon death or disability, and different
provisions for Cook. The 2022 form's Section 8 instead points to the vesting
date. An applicable executed agreement, amendment or further disclosure is
needed to establish which timing governs a particular executive. This is an
unresolved applicability and timing distinction between the inspected sources,
not an established contradiction for one identified award
\citep[printed p.~53]{apple2025proxy}.

Dividend-equivalent rights add a cash consequence to the same conditional
award. The filed form credits ordinary cash dividends using the outstanding
target units and pays the vested proportion in cash, subject to withholding
and the other terms. Suppose one dividend of \$0.25 occurs while all 1,000
target units are outstanding before severance, and there are no other
dividends or adjustments. The initial credit is \$250. A 1.5 performance
multiplier with full eligibility pays \$375 gross; under the assumed $3/5$
eligible fraction, the vested proportion is $900/1{,}000$, so \$225 is payable.
Applying the negative-return cap reduces that conditional payment to \$150.
This cash right is different from the hypothetical dividend reinvestment used
to measure shareholder return; adding it to that return would count a
different entitlement \citep[Exhibit 10.3, Section 5(b)]{apple2022plan}.

The model consequently needs distinct stages: determine the applicable award
and observations; calculate the return and rank under evidenced conventions;
apply the schedule and cap; determine surviving rights; and apply settlement
and withholding terms. The checked examples cover selected consequences with
explicit inputs. They do not recover the October 2026 outcome, reproduce all
individual awards, value compensation expense or establish diluted EPS. Those
questions require further data and their own accounting methods.
